\documentclass[11pt]{article}
\usepackage[margin=0.9in]{geometry}
\usepackage[T1]{fontenc}
\usepackage{lmodern}
\usepackage{microtype}
\usepackage{amsmath,amssymb}
\usepackage{booktabs,array,longtable}
\usepackage{enumitem}
\usepackage[hidelinks]{hyperref}
\usepackage{natbib}
\setlist{nosep,leftmargin=1.4em}
\setlength{\parindent}{0pt}
\setlength{\parskip}{0.5em}
\newcommand{\sz}{s_0}

\title{Our assumptions: origin, referee risk, and the S1--S3 simulations}
\author{SMC-BVS Writing Department}
\date{25 November 2026}

\begin{document}
\maketitle

\section*{How to read this}
Two reviewers read our assumptions A1--A10 and the computing conventions C1--C3 independently: Fable 5.2 and Astra
(gpt-6.2-astra). This note merges their answers. Where they disagree, both views are given. ``SAM'' is
\citet{SongLiang2015}; page numbers are journal pages, ``supp.'' means their supplement.

The assumptions describe a sequence of problems as $n$ grows. ``Holds'' for a design means the design, continued in $n$
with its structure fixed, satisfies the assumption. Where an assumption has explicit constants that our actual $(n,p,\sz)$
violate, we say so separately.

The designs, fixed by the owner on 24 November:
\begin{itemize}
  \item \textbf{S1}: AR(1) predictors, $\rho=0.6$; $n\in\{500,1000\}$, $p\in\{500,5000\}$, $\sz=10$;
    $\mathrm{SNR}\in\{0.5,1,2,3\}$, $\sigma=1$, and
    \[\beta^*=\mathrm{SNR}\sqrt{\log p/n}\,(2,-3,2,2,-3,3,-2,3,-2,3,0,\dots,0).\]
  \item \textbf{S2}: AR(1), $\rho\in\{0,0.35,0.7\}$; $(n,p,s)\in\{(50,1000,5),(100,1000,10)\}$; coefficients equal
    to $1$ (beta-types 1, 2); $\sigma^2=\beta^\top\Sigma\beta/\nu$ over ten values of $\nu$ from $0.05$ to $6$.
  \item \textbf{S3}: equicorrelated predictors, $\rho\in\{0.5,0.8\}$; $\sz\in\{5,10\}$; S1's $(n,p)$. The coefficients,
    their signs and $\sigma$ are not fixed in any source yet, so some S3 answers are conditional.
\end{itemize}

\section*{The answers at a glance}
Pushback: L low, M medium, H high. ``\,/\,'' separates the two reviewers (Fable\,/\,Astra) where they differ.

{\small
\begin{longtable}{@{}p{3.0cm}p{3.1cm}p{1.45cm}p{2.25cm}p{2.25cm}p{2.25cm}@{}}
\toprule
Assumption & Origin & Pushback & S1 & S2 & S3 \\
\midrule
\endhead
A1 Gaussian linear model & SAM, same & L & holds & holds & holds if $y$ is Gaussian \\
A2 $p\le n^A$ & not SAM (SAM allows faster $p$) & L\,/\,M & holds & holds & holds \\
A3 fixed $\sz$ & SAM, same & L--M\,/\,M & holds & holds; $n/\sz=10$ only & holds \\
A4(a) sub-Gaussian rows & not SAM; stronger than SAM's Cond.~1 & L & holds & holds & holds \\
A4(b) Gaussian rows & not SAM; stronger than anything SAM assumes & M--H & holds & holds & holds \\
A5(a) fixed raw signals & implicit in SAM's proofs & M\,/\,H & \textbf{fails} & holds; most cells weak & depends on $\beta^*$ \\
A5(b) fixed standardised signals & ours & H & \textbf{fails} & fails & fails \\
A6 eigenvalue floor & not SAM; standard & L\,/\,M & holds & holds & holds \\
A7 setting of Thm 2 & SAM (after Jiang 2007) & L\,/\,M & holds & holds; weak at $n=50$ & plausibly holds \\
A8 visibility & partial faithfulness; stronger than SAM's Cond.~5 & H & \textbf{fails} & holds & holds only if all signs agree \\
A9 small leakage set & ours & H & asymptotically only & asymptotically only & \textbf{fails} \\
A10 screening conditions & ours as a package & M\,/\,H & \textbf{fails} & \textbf{fails} & \textbf{fails} \\
C1--C3 idealised computer & not statistical & L\,/\,L--H & n/a & n/a & n/a \\
\bottomrule
\end{longtable}}

\section*{Assumption by assumption}

\textbf{A1. Gaussian linear model.} Same as SAM's model (their (1)--(2), pp.~947, 950). Every paper in this area
assumes it. A referee may ask about non-Gaussian errors; the planned misspecification result is our answer. It holds in
S1 and S2; in S3 once $y$ is generated with Gaussian noise.

\textbf{A2. At most polynomially many predictors, $p\le n^A$.} SAM allows more: $p$ up to $\exp(Cn^\alpha)$ for the
density result (Cor.~1, p.~954) and $n^{\kappa\log n}$ for selection (Cor.~S1, supp.~p.~16). Polynomial $p$ is the
assumption of \citet{BuhlmannKalischMaathuis2010}. Fable expects little pushback because $A$ can be any fixed number.
Astra expects some: a referee may say a fixed power excludes the ultrahigh-dimensional growth SAM covers. It holds in all
designs ($A\le1.77$).

\textbf{A3. A fixed number of true predictors.} SAM's strict sparsity, (4) on p.~950, used by all its selection
results. A referee who knows \citet{YangWainwrightJordan2016} will note that they let $\sz$ grow. It holds in all
designs, but in S2 $n/\sz=10$, so the limit statements say little there.

\textbf{A4. The law of the rows.} (a) asks every sparse combination of predictors to have Gaussian-like tails; (b) asks
exactly Gaussian rows, for Theorem~3. Neither is from SAM. (a) is the usual random-design condition
\citep[e.g.][]{NarisettyHe2014}. For centred predictors (a) implies SAM's Condition~1 (p.~953). The reviewers differ on
scope: Fable shows this when SAM's size cap is within our set size; Astra shows it for every cap (by convexity of
$t\mapsto e^{t^2}$), which would make the qualifier in our short version unnecessary. SAM's weaker Condition~2 (p.~954,
second moments only) is not in our short version's comparison and should be. (b) is where referees will push: SAM does
not need Gaussian rows, and (b) excludes genotypes and binary data, where the method is aimed. All three designs have
Gaussian rows, so both hold.

\textbf{A5. Fixed signals.} (a) asks the true coefficients to be fixed nonzero numbers on the raw scale. SAM states no
lower bound, but its proofs use fixed coefficients (supp.~pp.~5--6, 15), so (a) is SAM's in substance. The wider
literature asks only a bound that shrinks like $\sigma\sqrt{\log p/n}$ \citep[Cond.~3 of][]{FanLv2008}; a referee will
ask for that. (b) asks the coefficients to be fixed on the \emph{sample-standardised} scale; it is ours, and no data are
generated that way. Both reviewers find:
\begin{itemize}
  \item S1 violates (a) and (b) by construction: $\beta^*\propto\sqrt{\log p/n}$, the smallest coefficient runs from
    $0.08$ to $0.78$ over the cells and shrinks with $n$. Our short version's sentence ``it holds \dots as in our
    simulations'' is no longer true.
  \item S2 meets (a) in the letter (coefficients $1$), but most $\nu$ cells are weak in noise units at $n=50,100$.
  \item (b) fails in every design that draws $y$ from raw coefficients, because the standardised coefficient is the raw
    one times a random sample standard deviation.
\end{itemize}

\textbf{A6. Eigenvalue floor.} No set of up to about $5(A+1)(\sz+1)\log n$ predictors may be nearly collinear in the
population. Standard sparse-eigenvalue condition. SAM's counterparts are Condition~3 (p.~955) and the Gram bounds
(22)--(23) of Theorem~S (supp.~p.~13), whose floor may vanish; ours is fixed but has no ceiling. Fable: low pushback.
Astra: medium, because it also forbids near-duplicate null predictors that do not affect the truth. It holds in all
designs: the floors are $0.25$ (S1), $1$, $0.48$, $0.18$ (S2), and $0.5$, $0.2$ (S3). Equicorrelation $0.8$ does not
break it.

\textbf{A7. The setting of Theorem~2.} SAM's density-consistency setting, which follows Jiang (2007): our Theorem~2 has
the form of their (15), p.~952, with a population floor $c_\psi/n$ in place of their ridge prior. No beta-min. It is the
one assumption set that holds in every design. At $n=50$ the conclusion is weak (radius at least $0.28$ on a distance at
most $1$), and the constants $M_0,n_0$ have no numerical value. Our short version's remark that SAM's density result
assumes ``nothing about collinearity'' is true only for SAM's ridge prior.

\textbf{A8. Visibility.} Every true predictor keeps a covariance with the signal bounded away from $0$ after
conditioning on any $15$ other predictors. SAM's Condition~5 (p.~958) is the marginal case with no margin, and serves
only its screening step. Fable found published precedents that our documents miss: partial faithfulness with a margin
\citep[(D4) of][]{BuhlmannKalischMaathuis2010}, the marginal condition of \citet{FanLv2008}, and the conditional one of
\citet{BarutFanVerhasselt2016}. Our earlier ``no precedent found'' is wrong. Expect high pushback: Bühlmann et al.\ call
their own version restrictive, and it fails by cancellation when signs are mixed and predictors correlated.
\begin{itemize}
  \item S1 fails: with S1's sign pattern the smallest conditional covariance is $0.2\%$ of the coefficient, predictor~5
    has the wrong marginal sign, and the margin shrinks with $n$. The worst conditioning sets are not blocks the method
    would form, so this says A8 is too strong, not that the method fails.
  \item S2 holds: all coefficients are positive and correlations nonnegative.
  \item S3 holds if all signs agree and fails or nearly fails with mixed signs (exactly zero at $\rho=0.5$ for alternating
    $\pm1$). \citet{ShinBhattacharyaJohnson2018}, cited for S3, use random signs.
\end{itemize}

\textbf{A9. Small leakage set.} Only a set $Q$ of fewer than $\sqrt n$ predictors may imitate the signal, the rest must
be explainable locally, and the constants are tight ($15b_H+b<\tfrac12$). It is ours; SAM has only an informal remark
(p.~958). High pushback: long, uncheckable on data, written for our own search. Both reviewers find it can hold only as
a limit: at our sample sizes the size conditions fail in every cell (for instance the model-size cap needs $K\ge341$ when
$K=250$ at $n=500$; $|H_l|\le n^{b_H}$ allows at most $1.26$ predictors at $n=1000$). In S1 and S2 it plausibly holds as
$n\to\infty$. In S3 it fails outright: every null predictor stays correlated with the signal after conditioning, so $Q$
would have to contain all $p$ predictors.

\textbf{A10. Conditions for certified screening.} A floor on the observed design, a lower and an $\ell_1$ bound on the
coefficients, and a sample-size inequality with the constant $1/200$. The parts are standard; the package is ours.
Condition (c) fails in every cell, even with zero budget, by three or more orders of magnitude: its right side is at
most $n/(200\sz^2)$, about $0.05$ in S1, against a left side of at least $12$. S1 would need $n$ of about $3\times10^6$.
Fable rates pushback medium (the form is fine, the constant is the problem); Astra high (a referee may call Theorem~4
empty at practical sizes).

\textbf{C1--C3. The idealised computer.} Exact arithmetic, ideal random numbers, a machine that does not fail. Not
statistical and not checkable from a design. Every theorem about an algorithm idealises this way; SAM's sampler
statements are limits in the number of iterations (supp.~p.~13). Proposition~5 (unbiased evidence) is the classical SMC
identity under C3. Fable: low pushback. Astra: high if Theorem~3's delivery clause is read as a promise about real
hardware, and it notes that delivery for every $n$ needs the stronger convention with lifted integer limits, not C2
alone.

\section*{Verdict}
\textbf{Standard.} A1 and A3 are SAM's. A5(a) is SAM's in substance. A7 is SAM's density setting with one extra mild
condition. A2, A4(a), A6 and C1--C3 are standard, though not from SAM.

\textbf{Ours.} A5(b), A9 and the constant in A10(c). A8 in this form, though it should be cited as partial
faithfulness.

\textbf{At risk.} By theorem, under S1--S3:
\begin{itemize}
  \item \textbf{Theorem 3} (the reported model is the true model) covers none of the three designs. S1 fails A5(b) and
    A8; S2 fails A5(b) and A9's constants; S3 fails A9, and A8 unless signs agree. A5(b), A8 and A9 also carry the
    highest referee risk.
  \item \textbf{Theorem 1} (posterior concentration) does not cover S1, whose signals shrink like $\sqrt{\log p/n}$. It
    covers S2 in the letter and S3 if its coefficients are fixed numbers.
  \item \textbf{Theorem 4} (certified screening) is empty at our sizes: A10(c) fails everywhere.
  \item \textbf{Theorem 2} (fitted distribution) and \textbf{Proposition 5} (unbiased evidence) cover all three designs.
\end{itemize}

\textbf{For the owner to decide.} Either add one design that Theorem~3 covers (for example AR(1) with fixed
coefficients of one sign), or say plainly that S1--S3 test the method outside Theorem~3, with Theorem~2 as the result
that covers them. S3's coefficients, signs and $\sigma$ also need to be fixed; with mixed signs S3 fails A8.

\section*{Where the reviewers disagree}
\begin{itemize}
  \item Pushback levels: A2, A3, A6, A7 (Fable low, Astra medium); A5(a), A10 (Fable medium, Astra high); C1--C3
    (Fable low, Astra high if read as a hardware promise).
  \item A4(a) versus SAM's Condition~1: Fable shows the implication when SAM's cap is within our set size; Astra shows it
    for every cap. Astra's argument is short and looks right; if confirmed, drop the qualifier in the short version.
  \item A8 origin: only Fable searched outside SAM and found the partial-faithfulness precedents.
\end{itemize}

\section*{Corrections to our earlier documents}
\begin{enumerate}
  \item Short version, A5: ``as in our simulations'' is false for S1.
  \item Short version, SAM list: add SAM's Condition~2 (p.~954) and Theorem~S with (22)--(23), (26) (supp.~pp.~13,
    15), the SAM result that matches our Theorem~1.
  \item Short version, A7 row: ``nothing about collinearity'' holds only for SAM's ridge prior.
  \item A8: cite \citet{BuhlmannKalischMaathuis2010}, \citet{FanLv2008}, \citet{BarutFanVerhasselt2016}; ``no precedent''
    is wrong.
  \item Short version, ``Reasonable?'' column: rows A6, A9, A10 still describe the old simulation suite.
\end{enumerate}

\bibliographystyle{abbrvnat}
\bibliography{../../shared/refs}
\end{document}
